\section{Related Work}\label{sec:related}

\paragraph{RL for business process management.} \todo{resource allocation with deep RL,
prescriptive process monitoring, next-best-action recommendation} \aepn{}s and
GymPN~\cite{aepn2023,aepn_repr2025,gympn2025} provide the modelling framework and the
environment we build on.

\paragraph{Structured credit assignment.} Factored policy gradients exploit a given
influence structure between action factors and reward factors~\cite{spooner2021fpg};
action-dependent factorised baselines reduce policy-gradient variance~\cite{wu2018adfb};
value decomposition~\cite{vdn2018}, counterfactual baselines~\cite{coma2018} and
difference rewards~\cite{difference2002} address credit in cooperative multi-agent
settings. We derive the influence structure from the semantics of the process model and
state the scheduling conditions (local actor, component-scoped postponement) under which
it is valid.

\paragraph{Non-work-conserving control.} In queueing systems with heterogeneous servers it
can be optimal to leave a slow server idle~\cite{linkumar1984}. Component-scoped
postponement makes such policies learnable per process. \todo{more queueing-control
references}

\paragraph{Graph neural networks for decision making.} Learned dispatching policies over
graph representations have been studied in combinatorial
optimisation~\cite{zhang2020jobshop}. Relational graph networks~\cite{rgcn2018,rgat2019,hu2020hgt}
underlie our policy network.
